\chapter*{Introduction g\'{e}n\'{e}rale}
\addcontentsline{toc}{chapter}{Introduction g\'{e}n\'{e}rale}

\indent Dans un contexte \'{e}conomique en constante \'{e}volution, la gestion des demandes clients dans le secteur des pi\`{e}ces de rechange automobile constitue un d\'{e}fi majeur pour les entreprises. En effet, un processus de traitement efficace joue un r\^{o}le d\'{e}terminant dans la satisfaction client et la performance op\'{e}rationnelle \`{a} long terme \cite{microservices}. Selon des \'{e}tudes r\'{e}centes dans le domaine de la relation client \cite{rag}, environ 35\% du temps des op\'{e}rateurs est consacr\'{e} \`{a} des t\^{a}ches r\'{e}p\'{e}titives de saisie et de recherche d'information. Par ailleurs, 42\% des entreprises d\'{e}clarent avoir des difficult\'{e}s \`{a} standardiser leurs d\'{e}cisions de garantie, et 58\% se sentent submerg\'{e}es par le volume de dossiers \`{a} traiter \cite{whisper}.

\indent De nombreuses entreprises ont cherch\'{e} \`{a} am\'{e}liorer ce processus en adoptant des outils num\'{e}riques. Plusieurs solutions logicielles et plateformes existent d\'{e}j\`{a} pour g\'{e}rer les demandes clients, telles que Zendesk, Freshdesk ou Salesforce. Cependant, bien que ces solutions soient performantes, elles ne r\'{e}pondent pas toujours aux besoins sp\'{e}cifiques du secteur des pi\`{e}ces de rechange automobile. De plus, certaines de ces plateformes manquent de flexibilit\'{e} pour s'adapter aux processus m\'{e}tier sp\'{e}cifiques \`{a} chaque entreprise.

\indent Chez Insy2s, le processus actuel de traitement des demandes est supervis\'{e} par les chefs d'agence et les op\'{e}rateurs. Cependant, l'actuel circuit de gestion pr\'{e}sente plusieurs difficult\'{e}s. En effet, le suivi manuel et non outill\'{e} des diff\'{e}rentes \'{e}tapes s'av\`{e}re particuli\`{e}rement chronophage et exigeant pour les op\'{e}rateurs. Par ailleurs, les d\'{e}cisions de garantie sont prises de mani\`{e}re subjective, sans standardisation ni tra\c{c}abilit\'{e}. Cette situation g\'{e}n\`{e}re des incoh\'{e}rences dans le traitement des dossiers, ce qui complique la gestion et ralentit la productivit\'{e}.

\indent C'est dans ce contexte que la soci\'{e}t\'{e} Insy2s, sp\'{e}cialis\'{e}e dans le d\'{e}veloppement d'applications et les services d'ing\'{e}nierie informatique, vise \`{a} am\'{e}liorer l'exp\'{e}rience de gestion des demandes clients dans le secteur automobile.

\indent L'objectif de ce projet est de d\'{e}velopper une solution num\'{e}rique permettant de centraliser les informations, de simplifier la gestion des diff\'{e}rentes \'{e}tapes de traitement et d'am\'{e}liorer le processus de prise de d\'{e}cision gr\^{a}ce \`{a} l'intelligence artificielle. La plateforme propos\'{e}e, baptis\'{e}e \textbf{AutoAssist}, repose sur une architecture microservices compos\'{e}e de onze services ind\'{e}pendants, int\'{e}grant des technologies d'IA pour la transcription vocale (Whisper), l'analyse du langage naturel (spaCy) et l'aide \`{a} la d\'{e}cision (Gemma~4 via RAG).

\indent Pour assurer la r\'{e}ussite de ce projet, nous avons adopt\'{e} la m\'{e}thodologie agile \textbf{Scrum}, qui favorise l'adaptabilit\'{e}, la collaboration et l'am\'{e}lioration continue tout au long du d\'{e}veloppement.

\indent Ce rapport est structur\'{e} en trois chapitres, chacun correspondant \`{a} une \'{e}tape cl\'{e} du d\'{e}veloppement de notre projet.

\indent Le premier chapitre, intitul\'{e} \og Analyse et sp\'{e}cification des besoins \fg{}, pr\'{e}sente tout d'abord l'analyse de l'existant. Il expose ensuite les besoins du projet en d\'{e}taillant les acteurs concern\'{e}s, ainsi que les besoins fonctionnels et non fonctionnels. Ce chapitre inclut \'{e}galement la m\'{e}thodologie de travail adopt\'{e}e pour mener \`{a} bien ce projet.

\indent Le deuxi\`{e}me chapitre, intitul\'{e} \og Conception \fg{}, est consacr\'{e} \`{a} la mod\'{e}lisation de la solution. \`{A} travers des diagrammes UML et des descriptions architecturales, il illustre la structure du syst\`{e}me et les interactions entre ses diff\'{e}rents composants.

\indent Le troisi\`{e}me chapitre, intitul\'{e} \og R\'{e}alisation \fg{}, \'{e}num\`{e}re les diff\'{e}rents outils mat\'{e}riels et logiciels utilis\'{e}s. Il d\'{e}crit la mise en \oe uvre concr\`{e}te du projet ainsi que les r\'{e}sultats obtenus, illustr\'{e}s par des captures d'\'{e}cran.

\indent Ce rapport se cl\^{o}ture par une conclusion g\'{e}n\'{e}rale et perspectives qui synth\'{e}tise le travail r\'{e}alis\'{e} au cours de ce projet, pr\'{e}sente ses apports et projette les \'{e}ventuelles perspectives.
